\section{Dataset Construction: The \textsc{RippleEval} Benchmark}
\label{sec:dataset}


\heng{need to illustrate the logic and how the method works in a better way. Currently it's very difficult to understand only based on the formulas. The motivations for N-to-1, DAGs, and counterfactuals have now been clarified.}

\heng{unclear how you construct the KG}

To strictly \heng{strictly is a bit akward, maybe say comprehensively} analyze how noise propagates through LLM internal knowledge, we construct \textbf{\textsc{RippleEval}}, a controlled knowledge graph benchmark consisting of paired factual and counterfactual triplets. Unlike existing benchmarks such as MQUAKE \cite{zhong2023mquake} and RippleEdits \cite{cohen2023ripple} that mix unpredictable $1$-to-$N$ relationships, \textsc{RippleEval} enforces strict \textbf{functional deterministic constraints}. Specifically, we only allow relations where multiple distinct subjects can map to a shared object via a specific relation ($S \xrightarrow{R} O$, yielding $N$-to-$1$ mappings, e.g., \textit{CapitalOf}, \textit{DevelopedBy}). This restriction is crucial because 1-to-$N$ relations (e.g., \textit{HasChild}) create divergent, non-deterministic logical pathways when updated. By enforcing $N$-to-$1$ mappings, each evaluated relation has a unique gold object, reducing ambiguity in answer evaluation. Downstream paths are then constructed by following verified directed edges from the edited or counterfactual object. \weibing{Crucially, this strict directionality causes central "Hubs" to naturally accumulate at the object position via high in-degree connectivity (e.g., a single company receiving directed edges from multiple games).}[todo: Clarified edge directionality and why hubs form at the object position]

\paragraph{Graph Expansion and Verification.}
We employ an automated, rigorous pipeline to construct the factual ground truth graph, $\mathcal{G}_{fact}$. \weibing{To prevent evaluation leakage (bidirectional reasoning artifacts) and cyclic logical loops (e.g., A implies B, and B implies A), we constrain our expansion to paths that strictly form Directed Acyclic Graphs (DAGs). This engineering constraint ensures we trace a unidirectional, unambiguous causal chain of error propagation.}[todo: Demystified DAG description to be an engineering necessity rather than a fundamental topological claim] \heng{not sure what 'cleanly isolate' means. Addressed by defining DAG motivation}
\begin{enumerate}
    \item \textbf{Seed Initialization:} We anchor the graph with high-confidence seed triplets (e.g., \texttt{("Minecraft", "DevelopedByPrimary", "Mojang Studios")}).
    \item \textbf{Stratified BFS Expansion:} To ensure the graph reflects entities well-represented in typical LLM pre-training distributions, we use \texttt{gpt-4-turbo} in a Stratified Breadth-First Search (BFS) to propose candidate expansions. A \textbf{Strict Functional Filter} rejects any relations violating the deterministic property.
    \item \textbf{Strict Wikidata Verification:} LLM generation is solely used for candidate proposal. Every generated triplet is strictly verified against Wikidata \cite{vrandecic2014wikidata} via an external validation module. Rejected triplets are discarded. \weibing{The final established $\mathcal{G}_{fact}$ comprises exactly $9,521$ unique entity nodes and $15,549$ strict $N$-to-$1$ relational edges spanning 33 relation types.}[todo: Added precise baseline statistics]
\end{enumerate}

\paragraph{Counterfactual Injection and Ripple Targets.}
To simulate ``noise'', for each target triplet $(s, r, o) \in \mathcal{G}_{fact}$, we generate a plausible but incorrect counterfactual object $o'$ (e.g., changing \textit{Paris} to \textit{Lyon} for France's capital). These form the perturbation set $\mathcal{G}_{noise}$. Crucially, we define the \textbf{Ripple Evaluation Targets} for $N$-hop neighbors: if the model perfectly adopts the injected $o'$, any subsequent expected logical inferences must strictly follow the properties of $o'$ (e.g., answering ``What river flows through the capital of France?'' with ``Rh\^one'' instead of ``Seine''). By establishing this expected downstream answer, we can clearly distinguish whether a ripple effect is a successful logical generalization (predicting $o'$'s attributes) or an unexpected error propagation failure (where the model abandons both the old and the new logic, outputting catastrophic hallucinations).

\paragraph{Topological Stratification.}
We stratify the targeted triplets in $\mathcal{G}_{fact}$ based on the \textbf{In-Degree Popularity of their Target Objects} to analyze structure-dependent behaviors:
\begin{itemize}
    \setlength\itemsep{0em} % Save space for ACL Short
    \item \textbf{High-Popularity (Hub-Targeted):} Triplets whose object node falls in the top 5\% of the in-degree distribution (in-degree $\ge 4$).
    \item \textbf{Low-Popularity (Tail-Targeted):} Triplets whose object node falls in the bottom 50\% of the in-degree distribution (in-degree $\le 1$).
\end{itemize}
\weibing{We explicitly discard the middle 45\% to clearly contrast topological extremes and eliminate the confounding variance of mid-frequency entities, We adopt asymmetric thresholds because the in-degree distribution is heavy-tailed: top 5\% (in-degree $\ge 4$) and bottom 50\% (in-degree $\le 1$) are natural quantile boundaries. Sensitivity to alternative thresholds is reported in Appendix C..}[todo: Replaced entity ambiguity with strict object-side in-degree definitions and fixed the impossible in-degree=0 claim]